%%%%%%%%%%%%%%%%%%%%%%%%%%%%%%%%%%%%%%%%%%%%%%%%%%%%%%%%%%%%%%%%%%
%%%
%%%                   report_template_LaTeX
%%%                          main.tex
%%%   <https://github.com/matsukawa-group/report_template_LaTeX>
%%%
%%%%%%%%%%%%%%%%%%%%%%%%%%%%%%%%%%%%%%%%%%%%%%%%%%%%%%%%%%%%%%%%%%

%%% 文書クラスの設定 %%%
\documentclass[
    paper=a4paper,      % A4 用紙サイズ
    article,            % article 相当の文書クラス
    fontsize=11pt,      % 欧文サイズ 11 pt
    jafontsize=11pt,    % 和文サイズ 11 pt
    head_space=25mm,    % 天の余白
    foot_space=25mm,    % 地の余白
    gutter=20mm,        % のどの余白
    fore-edge=20mm      % 小口の余白
    ]{jlreq}            % jlreq クラスを使用

%==================================================
% パッケージの読み込み
%==================================================

\usepackage{settings}

%==================================================
% 基本設定
%==================================================

% 文書タイトル
\newcommand{\DocTitle}{%
  ここに文書のタイトルを入れます
}

% 日付
\newcommand{\DocDate}{%
  2026 年 6 月 1 日%
}

% 表紙を付けるかどうか
\newif\ifcover
\covertrue
% \coverfalse

%==================================================
% 著者の登録
%==================================================

\addauthor
  {著者名}
  {著者の所属・学年}
  {author@example.com}

% 第二著者を追加する場合
%
% \addauthor
%   {第二著者}
%   {第二著者の所属・学年}
%   {second@example.com}


%==================================================
% 本文
%==================================================

\begin{document}

% \maketitle
\makedocumenttitle

\thispagestyle{cover}

% 目次が不要な場合はコメントアウト
\tableofcontents
\clearpage

\section{これは見出し}
\label{sec:heading}

\begin{lstlisting}[language=TeX]
\section{これは見出し}
\label{sec:heading}
\end{lstlisting}

\subsection{これは小見出し}
\label{ssec:subheading}

\begin{lstlisting}[language=TeX]
\subsection{これは小見出し}
\label{ssec:subheading}
\end{lstlisting}

\subsubsection{これはさらに小さい見出し}
\label{sssec:subsubheading}

\begin{lstlisting}[language=TeX]
\subsubsection{これはさらに小さい見出し}
\label{sssec:subsubheading}
\end{lstlisting}

\section{\LaTeX 文書の基本的な書き方}

ここでは \LaTeX 文書の基本的な書き方について説明します．

\subsection{コメント文・段落の変え方・強調の仕方}
\label{ssec:paragraph}

% ← これはコメント文にするコマンドです．
% \jalipsum[1-10]{wagahai} これはダミーテキストです．
\jalipsum[1-2]{wagahai} % ↓ 空白行を入れると改段落されます．

\textbf{\color{red}{段落を変えたいときは空白行を入れます．}} % 強調のために色と太さを変えました．
\jalipsum[3-4]{wagahai}

\begin{lstlisting}[language=TeX]
% ← これはコメント文にするコマンドです．
% \jalipsum[1-10]{wagahai} これはダミーテキストです．
\jalipsum[1-2]{wagahai} % ↓ 空白行を入れると改段落されます．

\textbf{\color{red}{段落を変えたいときは空白行を入れます．}} % 強調のために色と太さを変えました．
\jalipsum[3-4]{wagahai}
\end{lstlisting}

\subsection{相互参照}
\label{ssec:reference}

\LaTeX の強みの一つは相互参照が簡単にできることです．
\LaTeX では \verb|\label{xxx}| のように \verb|xxx| という名前のラベルを定義できます．
例えば，この節は見出しのところで
\begin{lstlisting}[language=TeX]
\subsection{相互参照}
\label{ssec:reference}
\end{lstlisting}
としているので，この見出しを参照したいときは
\begin{lstlisting}[language=TeX]
第~\ref{ssec:reference}~節
\end{lstlisting}
とすれば「第~\ref{ssec:reference}~節」のように出力できます．
このような相互参照は見出しだけでなく，数式や図表などあらゆる要素に対して行うことができます．
\verb|xxx| の部分の文字列は自由に設定できるので，ジャンルごとにわかりやすいラベルをつけるといいでしょう\footnote{例：節の場合は \texttt{\textbackslash label\{sec:xxx\}}，小節の場合は \texttt{\textbackslash label\{ssec:xxx\}}，小々節の場合は \texttt{\textbackslash label\{sssec:xxx\}}，式の場合は \texttt{\textbackslash label\{eq:xxx\}}，図の場合は \texttt{\textbackslash label\{fig:xxx\}}，表の場合は \texttt{\textbackslash label\{tb:xxx\}} など，自分のわかりやすいラベルをつけるといいでしょう．ちなみにここに表示しているような脚注は \texttt{\textbackslash footnote\{脚注内の文字列\}} で出力します．}．
また，コード中に挿入している \verb|~|（チルダ）を使用することで，この箇所での改行を防ぐことができます．

\subsection{強制改ページ}
\label{ssec:pagebreak}

紙面の都合で強制的に改ページしたいときは \verb|\clearpage| コマンドを使用します．
例えば，この段落の後で \verb|\clearpage| とすれば，この段落の後で改ページされます．

\clearpage

\subsection{箇条書き}
\label{ssec:list}

箇条書きには番号なしのものと番号付きのものがあります．

\subsubsection{番号なしの箇条書き}
\label{sssec:itemize}

番号なしの箇条書きの例として，明治大学のキャンパスと学部を挙げます．
箇条書きは入れ子にすることもできます．

\begin{multicols}{2}
  \begin{itemize}
    \item 和泉，駿河台キャンパス
    \begin{itemize}
      \item 法学部
      \item 商学部
      \item 政治経済学部
      \item 文学部
      \item 経営学部
      \item 情報コミュニケーション学部
    \end{itemize}
    \item 生田キャンパス
    \begin{itemize}
      \item 理工学部
      \begin{itemize}
        \item 機械工学科
        \item 機械情報工学科
      \end{itemize}
      \item 農学部
    \end{itemize}
    \item 中野キャンパス
    \begin{itemize}
      \item 国際日本学部
      \item 総合数理学部
    \end{itemize}
  \end{itemize}
  \columnbreak
\begin{lstlisting}[language=TeX]
\begin{itemize}
  \item 和泉，駿河台キャンパス
  \begin{itemize}
    \item 法学部
    \item 商学部
    \item 政治経済学部
    \item 文学部
    \item 経営学部
    \item 情報コミュニケーション学部
  \end{itemize}
  \item 生田キャンパス
  \begin{itemize}
    \item 理工学部
    \begin{itemize}
      \item 機械工学科
      \item 機械情報工学科
    \end{itemize}
    \item 農学部
  \end{itemize}
  \item 中野キャンパス
  \begin{itemize}
    \item 国際日本学部
    \item 総合数理学部
  \end{itemize}
\end{itemize}
\end{lstlisting}
\end{multicols}

\vspace{-2mm}

\subsubsection{番号付きの箇条書き}
\label{sssec:enumerate}

番号付きの箇条書きの例として，日本の苗字ランキングの上位を挙げます．
番号付きの箇条書きも入れ子にすることができます．

\begin{multicols}{2}
  \begin{enumerate}
    \item 佐藤
    \item 鈴木
    \item 高橋
    \begin{enumerate}
      \item 高橋
      \item 髙橋
    \end{enumerate}
    \item 田中
    \item 伊藤
  \end{enumerate}
  \columnbreak
\begin{lstlisting}[language=TeX]
\begin{enumerate}
  \item 佐藤
  \item 鈴木
  \item 高橋
  \begin{enumerate}
    \item 高橋
    \item 髙橋
  \end{enumerate}
  \item 田中
  \item 伊藤
\end{enumerate}
\end{lstlisting}
\end{multicols}

\subsection{数式}
\label{ssec:math}

次に \LaTeX における数式の書き方について説明します．

\subsubsection{基本的な数式の書き方}
\label{sssec:math-basic}

\LaTeX で文章中に数式を組み込むインライン数式の場合は，\verb|$| と \verb|$| で表示したい数式を囲って \verb|$E = mc^2$| とすると $E = mc^2$ のように出力できます．

最も基本的な別行立ての数式は
\begin{equation}
    \pdv{u_r}{t} + (\mathbf{u}\cdot\nabla)u_r - \frac{u_\theta^2}{r} = -\frac{1}{\rho}\pdv{p}{r} + \nu\ab(\nabla^2 u_r - \frac{u_r}{r^2} - \frac{2}{r^2}\pdv{u_\theta}{\theta})
    \label{eq:NSr}
\end{equation}
\begin{lstlisting}[language=TeX]
\begin{equation}
    \pdv{u_r}{t} + (\mathbf{u}\cdot\nabla)u_r - \frac{u_\theta^2}{r} = -\frac{1}{\rho}\pdv{p}{r} + \nu\ab(\nabla^2 u_r - \frac{u_r}{r^2} - \frac{2}{r^2}\pdv{u_\theta}{\theta})
    \label{eq:NSr}
\end{equation}
\end{lstlisting}
のように \verb|equation| 環境で記述します．
式~\eqref{eq:NSr} の数式は \LaTeX でサポートされている最も標準的なコマンドと \href{https://ctan.org/pkg/physics2}{\texttt{physics2}} パッケージ，\href{https://ctan.org/pkg/fixdif}{\texttt{fixdif}} パッケージ，\href{https://ctan.org/pkg/derivative}{\texttt{derivative}} パッケージで記述しています．
上付き添え字はキャレット \verb|^|，下付き添え字はアンダースコア \verb|_| を用いて表現します．
したがって，$u_\theta^2$ は \verb|u_\theta^2| と書きます．
ここで注意点として，添え字が $u_\theta^2$ のように一文字であれば問題無いのですが，$R_{ij}$ のように二文字以上の場合には $R_{ij}$ のように括弧 \verb|{}| で囲んでください．
分数は \verb|\frac{分子}{分母}| で書きますが，インライン数式で $1/2$ のようにスラッシュで分数表記したいときは \verb|1/2| のようにスラッシュを入力すれば大丈夫です．
表~\ref{tb:greek} のように $\rho$ や $\theta$ といったギリシャ文字も出力できるほか，$\nabla$ や $\sin$，$\log$ のような数学で使う関数の類もコマンドが存在しています（例：\verb|\nabla|, \verb|\sin|, \verb|\log|）．
$\sin$ や $\log$ は通常アップライト体（直立体，Roman 体）で書きます．
$sinx$ などと書くことのないように気をつけましょう．
また，数式中で \LaTeX コマンドを使わずに変換で出した全角のθを入れている人をときどき見かけるので気をつけましょう．

式~\eqref{eq:NSr} の左辺と右辺で丸括弧の大きさが異なることにも注目してください．
左辺は特に何もしていませんが，右辺は分数がある分，縦方向に括弧の長さが必要です．
括弧の大きさを自動で調整するコマンドとして，\verb|\ab(xxx)|コマンドがあります．


\begin{table}[t]
    \centering
    \caption{ギリシャ文字の出力方法．}
    \label{tb:greek}
    \begin{tabular}{c|c|c|c|c} \hline\hline
        英語名 & 大文字出力 & 大文字入力 & 小文字出力 & 小文字入力 \\ \hline
        alpha & $A$ & \verb|A| & $\alpha$ & \verb|\alpha| \\ \hline
        beta & $B$ & \verb|B| & $\beta$ & \verb|\beta| \\ \hline
        gamma & $\Gamma$, $\varGamma$ & \verb|\Gamma|, \verb|\varGamma| & $\gamma$ & \verb|\gamma| \\ \hline
        delta & $\Delta$, $\varDelta$ & \verb|\Delta|, \verb|\varDelta| & $\delta$ & \verb|\delta| \\ \hline
        epsilon & $E$ & \verb|E| & $\epsilon$, $\varepsilon$ & \verb|\epsilon|, \verb|\varepsilon| \\ \hline
        zeta & $Z$ & \verb|Z| & $\zeta$ & \verb|\zeta| \\ \hline
        eta & $H$ & \verb|H| & $\eta$ & \verb|\eta| \\ \hline
        theta & $\Theta$, $\varTheta$ & \verb|\Theta|, \verb|\varTheta| & $\theta$, $\vartheta$ & \verb|\theta|, \verb|\vartheta| \\ \hline
        iota & $I$ & \verb|I| & $\iota$ & \verb|\iota| \\ \hline
        kappa & $K$ & \verb|K| & $\kappa$, $\varkappa$ & \verb|\kappa|, \verb|\varkappa| \\ \hline
        lambda & $\Lambda$, $\varLambda$ & \verb|\Lambda|, \verb|\varLambda| & $\lambda$ & \verb|\lambda| \\ \hline
        mu & $M$ & \verb|M| & $\mu$ & \verb|\mu| \\ \hline
        nu & $N$ & \verb|N| & $\nu$ & \verb|\nu| \\ \hline
        xi & $\Xi$, $\varXi$ & \verb|\Xi|, \verb|\varXi| & $\xi$ & \verb|\xi| \\ \hline
        omicron & $O$ & \verb|O| & $o$ & \verb|o| \\ \hline
        pi & $\Pi$, $\varPi$ & \verb|\Pi|, \verb|\varPi| & $\pi$, $\varpi$ & \verb|\pi|, \verb|\varpi| \\ \hline
        rho & $P$ & \verb|P| & $\rho$, $\varrho$ & \verb|\rho|, \verb|\varrho| \\ \hline
        sigma & $\Sigma$, $\varSigma$ & \verb|\Sigma|, \verb|\varSigma| & $\sigma$, $\varsigma$ & \verb|\sigma|, \verb|\varsigma| \\ \hline
        tau & $T$ & \verb|T| & $\tau$ & \verb|\tau| \\ \hline
        upsilon & $\Upsilon$, $\varUpsilon$ & \verb|\Upsilon|, \verb|\varUpsilon| & $\upsilon$ & \verb|\upsilon| \\ \hline
        phi & $\Phi$, $\varPhi$ & \verb|\Phi|, \verb|\varPhi| & $\phi$, $\varphi$ & \verb|\phi|, \verb|\varphi| \\ \hline
        chi & $X$ & \verb|X| & $\chi$ & \verb|\chi| \\ \hline
        psi & $\Psi$, $\varPsi$ & \verb|\Psi|, \verb|\varPsi| & $\psi$ & \verb|\psi| \\ \hline
        omega & $\Omega$, $\varOmega$ & \verb|\Omega|, \verb|\varOmega| & $\omega$ & \verb|\omega| \\ \hline
    \end{tabular}
\end{table}



%%% 文献 %%%
% \biblist
% 使用する bst ファイル
% \bibliographystyle{jsme}
% 読み込む bib ファイル
% \bibliography{
%     mybib_en.bib,
%     mybib_jp.bib
% }

\end{document}
